\subsection{平坦模的张量积}

设 A、B 是两个环，E 是右 A-模，F 是 (A, B)-双模（《代数学》第二章 §1 第 14 目）。回顾（《代数学》第二章 §3 第 4 目）， \(E \otimes_{A} F\) 具有一个典范的右 B-模结构，对此结构

\[
(x \otimes y) b = x \otimes (y b) \quad \text {   for   } \quad x \in \mathrm{E}, y \in \mathrm{F}, b \in \mathrm{B}.
\]

命题 8. 设 A、B 是两个环，E 是右 A-模，F 是 (A, B)-双模。设 E 是平坦的，且 F 作为 B-模是平坦的。则 B-模 \(E \otimes_{A} F\) 是平坦的。

设 G 是左 B-模， \(G'\) 是 G 的子模。由于 F 作为右 B-模是平坦的，典范同态 \(F \otimes_{B} G' \to F \otimes_{B} G\) 是单射的。由于 E 是平坦的，典范同态

\[
\mathrm{E} \otimes_ {\mathbf {A}} (\mathrm{F} \otimes_ {\mathbf {B}} \mathrm{G} ^ {\prime}) \rightarrow \mathrm{E} \otimes_ {\mathbf {A}} (\mathrm{F} \otimes_ {\mathbf {B}} \mathrm{G})
\]

是单射的。由于 \(E \otimes_{A} (F \otimes_{B} G')\) 和 \(E \otimes_{A} (F \otimes_{B} G)\) 分别典范地等同于 \((E \otimes_{A} F) \otimes_{B} G'\) 和 \((E \otimes_{A} F) \otimes_{B} G\) （《代数学》第二章 §3 第 8 目命题 8），典范同态

\[
(\mathbf {E} \otimes_ {\mathbf {A}} \mathbf {F}) \otimes_ {\mathbf {B}} \mathbf {G} ^ {\prime} \rightarrow (\mathbf {E} \otimes_ {\mathbf {A}} \mathbf {F}) \otimes_ {\mathbf {B}} \mathbf {G}
\]

是单射的，这就证明了 \(E \otimes_{A} F\) 是一个平坦 B-模。

推论 1. 设 C 是交换环，E、F 是两个平坦 C-模。则 \(E \otimes_{C} F\) 是平坦 C-模。

F 是 (C, C)-双模，只需取 B = A = C 应用命题 8。

推论 2. 设 \(\rho\) 是环 A 到环 B 的同态。若 E 是平坦右 A-模，则通过把基环扩张到 B（《代数学》第二章 §5 第 1 目）得到的右 B-模 \(\rho^{*}(\mathrm{E}) = \mathrm{E}_{(\mathrm{B})}\) 是平坦的。

按定义 \(E_{(B)} = E \otimes_{A} B\) ，其中 B 经由 \(\rho\) 被视为 (A, B)-双模。由于右 B-模 \(B_{d}\) 是平坦的，只需应用命题 8。

推论 3. 设 R、S 是两个环， \(\phi: \mathbf{R} \to \mathbf{S}\) 是环同态。若 M 是平坦右 S-模，且 \(\phi_*(\mathbf{S}_d)\) 是平坦右 R-模，则 \(\phi_*(\mathbf{M})\) 是平坦右 R-模。

回顾 \(\phi_{*}(\mathbf{M})\) 是由 \(x.r = x.\phi (r)\) （对所有 \(x\in \mathbf{M}\) 和所有 \(r\in \mathbb{R}\) ）定义的右 R-模（《代数学》第二章 §1 第 13 目）。于是取 \(\mathrm{A} = \mathrm{S}\) 、 \(\mathrm{B} = \mathrm{R}\) 、 \(\mathrm{E} = \mathrm{M}\) 、 \(\mathrm{F} = \mathrm{S}\) 应用命题 8，其中 S 具有 \(\phi\) 所定义的 (S, R)-双模结构；这时右 R-模 M \(\otimes_{\mathbf{S}}\) S 恰好就是 \(\phi_{*}(\mathbf{M})\) 。

命题 9. 设 \((\mathbf{A}_{\alpha}, f_{\beta \alpha})\) 是环的一个正向系， \(\mathbf{A} = \varinjlim \mathbf{A}_{\alpha}\) 是它的正向极限， \((\mathbf{E}_{\alpha}, f_{\beta \alpha})\) 是以同一指标集为指标的一族右 \(\mathbf{A}_{\alpha}\) -模的正向系， \(\mathbf{E} = \varinjlim \mathbf{E}_{\alpha}\) 是它的正向极限，它是一个右 A-模（《代数学》第二章 §6 第 2 目）。若诸 \(\mathbf{E}_{\alpha}\) 中的每一个都是平坦 \(\mathbf{A}_{\alpha}\) -模，则 \(\mathbf{E}\) 是平坦 A-模。

设 \(E_{\alpha}^{\prime}=E_{\alpha}\otimes_{A_{\alpha}}A\) ，其中 A 被视为左 \(A_{\alpha}\) -模（经由典范同态 \(A_{\alpha}\to A\) ）；我们知道右 A-模 E 典范地同构于 \(\varinjlim E_{\alpha}^{\prime}\) （出处同上，命题 7 的推论 2）。由命题 8 的推论 2 得， \(E_{\alpha}^{\prime}\) 对所有 \(\alpha\) 都是平坦右 A-模，因此由第 3 目的命题 2，E 是平坦 A-模。

